\documentclass[12pt,a4paper]{article}
\usepackage[utf8]{inputenc}
\usepackage[russian]{babel}
\usepackage{multirow}

\begin{document}

\section*{Эксперименты с таблицами в LaTeX}

\subsection*{1. Простой пример таблицы}

\begin{tabular}{|l|c|r|}
\hline
Имя & Возраст & Оценка \\
\hline
Анна & 20 & 95 \\
Борис & 22 & 87 \\
Виктор & 21 & 92 \\
\hline
\end{tabular}

\vspace{1cm}

\subsection*{2. Разные выравнивания (l, c, r)}

\textbf{l - выравнивание по левому краю, c - по центру, r - по правому краю}

\vspace{0.5cm}

\begin{tabular}{|l|c|r|}
\hline
\textbf{Левое (l)} & \textbf{Центр (c)} & \textbf{Правое (r)} \\
\hline
Текст слева & Текст в центре & Текст справа \\
123 & 456 & 789 \\
Короткий & Средний текст & Длинный текст здесь \\
\hline
\end{tabular}

\vspace{1cm}

\subsection*{3. Слишком мало элементов в строке}

\textbf{Результат:} Пустые ячейки в конце строки

\vspace{0.5cm}

\begin{tabular}{|l|c|r|c|}
\hline
Столбец 1 & Столбец 2 & Столбец 3 & Столбец 4 \\
\hline
Полная & строка & со всеми & элементами \\
\hline
Только & два & & \\
\hline
Один & & & \\
\hline
\end{tabular}

\vspace{1cm}

\subsection*{4. Слишком много элементов в строке}

\textbf{Результат:} Ошибка компиляции! LaTeX выдаст ошибку:\\
\texttt{Extra alignment tab has been changed to \textbackslash cr}

\vspace{0.5cm}

% Закомментировано, чтобы избежать ошибки:
% \begin{tabular}{|l|c|r|}
% \hline
% Столбец 1 & Столбец 2 & Столбец 3 \\
% \hline
% Один & Два & Три & Четыре & Пять \\  % ОШИБКА!
% \hline
% \end{tabular}

\textit{(Пример закомментирован, так как вызывает ошибку компиляции)}

\vspace{1cm}

\subsection*{5. Использование multicolumn для объединения столбцов}

\textbf{Синтаксис:} \texttt{\textbackslash multicolumn\{число\_столбцов\}\{выравнивание\}\{текст\}}

\vspace{0.5cm}

\textbf{Пример 1: Заголовок на несколько столбцов}

\begin{tabular}{|l|c|r|c|}
\hline
\multicolumn{4}{|c|}{\textbf{Таблица успеваемости студентов}} \\
\hline
Имя & Математика & Физика & Средний балл \\
\hline
Алексей & 85 & 90 & 87.5 \\
Мария & 92 & 88 & 90.0 \\
Дмитрий & 78 & 82 & 80.0 \\
\hline
\multicolumn{3}{|r|}{\textbf{Средний балл группы:}} & \textbf{85.8} \\
\hline
\end{tabular}

\vspace{1cm}

\textbf{Пример 2: Объединение в середине таблицы}

\begin{tabular}{|l|c|c|c|}
\hline
Продукт & Январь & Февраль & Март \\
\hline
Яблоки & 100 & 120 & 110 \\
\hline
Груши & \multicolumn{2}{c|}{Нет данных} & 95 \\
\hline
Апельсины & 80 & 85 & 90 \\
\hline
\end{tabular}

\vspace{1cm}

\textbf{Пример 3: Сложная структура}

\begin{tabular}{|l|c|c|c|c|}
\hline
\multicolumn{5}{|c|}{\textbf{Отчет по продажам}} \\
\hline
\multirow{2}{*}{\textbf{Регион}} & \multicolumn{3}{c|}{\textbf{Кварталы}} & \multirow{2}{*}{\textbf{Итого}} \\
\cline{2-4}
 & Q1 & Q2 & Q3 & \\
\hline
Москва & 1000 & 1200 & 1100 & 3300 \\
\hline
Санкт-Петербург & 800 & 850 & 900 & 2550 \\
\hline
Регионы & 600 & 650 & 700 & 1950 \\
\hline
\multicolumn{4}{|r|}{\textbf{Всего:}} & \textbf{7800} \\
\hline
\end{tabular}

\vspace{1cm}

\subsection*{Выводы:}

\begin{itemize}
    \item \textbf{l, c, r} определяют выравнивание содержимого в столбцах
    \item Если элементов \textbf{меньше}, чем столбцов - остальные ячейки остаются пустыми
    \item Если элементов \textbf{больше}, чем столбцов - возникает ошибка компиляции
    \item \texttt{\textbackslash multicolumn} позволяет объединять столбцы и изменять их форматирование
    \item Можно комбинировать с \texttt{\textbackslash multirow} для объединения строк
\end{itemize}

\end{document}
